% !PRINTED 244
\weberpage{244}
If

\weq{3}{P = t^n + p_1 t^{n-1} + p_2 t^{n-2} + \cdots + p_{n-1} t + p_n}

denotes a function of degree $n$ irreducible with respect to the modulus $p$, in which for the sake of simplicity we assume the coefficient of the highest power $t^n$ to be equal to $1$, then from every other integral function $F(t)$ with integral coefficients a residue $\Phi(t)$ can be derived whose degree is lower than $n$, by setting (according to \S.\,3 of the first volume)

\weq{4}{F(t) = QP + \Phi(t)}

Here $\Phi(t)$ has integral coefficients, and we designate its relation to the function $F(t)$ as a congruence with respect to the modulus $P$.

\gesperrt{Accordingly, two functions $F(t)$, $F_1(t)$ that have the same residue $\Phi(t)$ are called congruent with respect to the modulus $P$}, which is expressed by the formula

\[
F(t) \equiv F_1(t)\ (\mathrm{mod}\ P)
\]

This is equivalent to saying that $F(t) - F_1(t)$ is divisible by $P$.

If two functions $F(t)$ and $F_1(t)$ do not give the same residue, but two residues that are congruent with respect to the modulus $p$, then the functions $F$, $F_1$ are called congruent with respect to the double modulus $P$, $p$, and one expresses this by a formula as follows:

\weq{5}{F(t) \equiv F_1(t)\ (\mathrm{mod}\ P, p).}

For this kind of congruence the following theorem holds:

\gesperrt{1. The product of two functions $F(t)$, $F_1(t)$ cannot be congruent to zero unless one of the factors is congruent to zero.}

The proof follows from the algorithm of the greatest common divisor. For if $P_1$ is an integral function of lower degree than $P$ which is not congruent to zero with respect to the modulus $p$, then one can determine a sequence of functions of the same kind $P_2$, $P_3$, $\ldots$ of decreasing degree, and the quotients $Q_1$, $Q_2$, $\ldots$ likewise as integral functions, in such a way that

\weq{6}{P \equiv Q_1 P_1 + P_2,\quad P_1 \equiv Q_2 P_2 + P_3, \ldots,\quad P_{\nu-2} \equiv Q_{\nu-1} P_{\nu-1} + P_\nu\ (\mathrm{mod}\ p)}

holds, and the sequence of these equations breaks off when $P_\nu$ has become a constant (an integer). This constant $P_\nu$
%CONT>>
